% ===================== CHAPITRE 5 — SPRINT 3 =========================
\chapter{Sprint 3 : Gestion des clients et des objectifs}

\section*{Introduction}
Ce sprint complète le référentiel commercial : la gestion du portefeuille clients,
leur affectation aux gestionnaires et le suivi des objectifs commerciaux.

\section{Objectifs du Sprint}
L'objectif est de permettre la \textbf{gestion des clients} (portefeuille), leur
\textbf{affectation à un gestionnaire} et le \textbf{suivi des objectifs
commerciaux} (ouvertures de comptes, crédits, épargne) par agence.

\section{Spécification des besoins}

\subsection{Diagramme de cas d'utilisation du Sprint 3}
La figure~\ref{fig:uc_s3} présente les cas d'utilisation du sprint.

\begin{figure}[htbp]
  \centering
  \includegraphics[width=0.82\textwidth]{uc_sprint3.png}
  \caption{Diagramme de cas d'utilisation du Sprint 3}
  \label{fig:uc_s3}
\end{figure}

\subsection{Sprint Backlog du sprint 3}
Le tableau~\ref{tab:bl_s3} détaille les user stories et les tâches techniques.

\begin{table}[htbp]
\centering
\caption{Sprint Backlog du sprint 3}
\label{tab:bl_s3}
\renewcommand{\arraystretch}{1.3}
\begin{tabularx}{\textwidth}{>{\bfseries}p{1cm}p{5.2cm}X}
\toprule
\textmd{\textbf{US}} & \textbf{User story} & \textbf{Tâches techniques}\\
\midrule
US6 & Gérer les clients (portefeuille). &
Développer le \texttt{ClientServlet} et le \texttt{ClientService} sur l'entité
\texttt{Client} ; page \texttt{client.jsp} ; gérer les clients BTK
(\texttt{BClientServlet}, \texttt{clients-btk.jsp}).\\
\midrule
US7 & Affecter un client à un gestionnaire. &
Associer chaque \texttt{Client} à un \texttt{Gestionnaire}
(\texttt{SK\_GESTIONNAIRE}) ; mettre à jour le portefeuille du gestionnaire.\\
\midrule
US8 & Suivre les objectifs commerciaux et la production. &
Développer l'\texttt{ObjectifServlet} et l'\texttt{ObjectifService} sur l'entité
\texttt{Objectif} ; pages \texttt{objectifs.jsp} et \texttt{objectifs-btk.jsp} ;
restituer comptes, crédits et épargne par agence.\\
\bottomrule
\end{tabularx}
\end{table}

\section{Analyse des besoins}

\subsection{Description textuelle des cas d'utilisation}

\cudesc{tab:cu_client}{Gérer les clients}
{Permettre à l'administrateur de gérer le portefeuille clients et d'affecter
chaque client à un gestionnaire.}
{Administrateur}
{L'administrateur doit être authentifié et le gestionnaire de rattachement doit
exister.}
{Le client est enregistré et associé à son gestionnaire dans la base.}
{(1) L'administrateur accède à l'interface de gestion des clients ; (2) il saisit
ou met à jour les informations du client ; (3) il sélectionne le gestionnaire
responsable ; (4) le \texttt{ClientService} enregistre le client et met à jour le
portefeuille.}

\cudesc{tab:cu_objectif}{Suivre les objectifs commerciaux}
{Permettre au directeur commercial de consulter les objectifs et la production par
agence.}
{Directeur commercial / Administrateur}
{L'acteur doit être authentifié et des objectifs doivent être renseignés pour la
période.}
{Les indicateurs de production (comptes, crédits, épargne) sont restitués par
agence.}
{(1) L'acteur accède à l'interface des objectifs ; (2) il sélectionne une agence
et une période ; (3) l'\texttt{ObjectifService} agrège la production ; (4) les
indicateurs sont affichés.}

\section{Conception}

\subsection{Diagramme de séquence : « Gérer un client »}
La figure~\ref{fig:seq_client} décrit l'enregistrement et l'affectation d'un
client.

\begin{figure}[htbp]
  \centering
  \includegraphics[width=0.92\textwidth]{seq_client.png}
  \caption{Diagramme de séquence -- Gestion d'un client}
  \label{fig:seq_client}
\end{figure}

Le \texttt{ClientServlet} transmet la demande au \texttt{ClientService}, qui
recherche le \texttt{Gestionnaire} concerné, associe le \texttt{Client} à ce
dernier, puis le persiste via l'\texttt{EntityManager}. Le portefeuille mis à jour
est renvoyé à la page \texttt{client.jsp}.

\section{Réalisation : interfaces utilisateur}
La figure~\ref{fig:ui_clients} illustre les interfaces de gestion des clients et
de suivi des objectifs. \emph{Les captures d'écran réelles seront insérées à cet
emplacement.}

\begin{figure}[htbp]
  \centering
  \fbox{\parbox[c][3.6cm][c]{0.8\textwidth}{\centering\itshape Capture à insérer
  -- gestion des clients et suivi des objectifs}}
  \caption{Interfaces de gestion des clients et des objectifs}
  \label{fig:ui_clients}
\end{figure}

\section*{Conclusion}
Ce sprint a complété la gestion commerciale (clients et objectifs). Le sprint
suivant aborde le pointage des employés, les workflows de demandes et les
notifications.
